\documentclass[11pt, a4paper]{article}

\usepackage[utf8]{inputenc}
\usepackage[T1]{fontenc}
\usepackage[spanish]{babel}
\usepackage{lmodern}
\usepackage[margin=2.2cm]{geometry}
\usepackage{xcolor}
\usepackage{fancyhdr}
\usepackage{enumitem}
\usepackage{titlesec}
\usepackage{tcolorbox}
\usepackage{amssymb}
\usepackage{parskip}
\usepackage{hyperref}
\usepackage{needspace}
\tcbuselibrary{skins, breakable}

% ---- Colores estilo data.lung ----
\definecolor{verdeTerm}{HTML}{1C8A3F}
\definecolor{verdeOsc}{HTML}{146A2F}
\definecolor{tinta}{HTML}{0B0B0B}
\definecolor{tinta2}{HTML}{4A4948}
\definecolor{mute}{HTML}{8F8D87}
\definecolor{plano}{HTML}{F7F4EB}
\definecolor{planoOsc}{HTML}{EFEBDD}
\definecolor{linea}{HTML}{DAD6C8}
\definecolor{ambar}{HTML}{C78A1E}

\hypersetup{
  colorlinks=true, linkcolor=verdeTerm, urlcolor=verdeTerm,
  pdftitle={Guia de estudio - Defensa TFM data.lung},
  pdfauthor={Daniel Tapia Diez}
}

% ---- Titulos con filete verde ----
\titleformat{\section}
  {\Large\bfseries\color{tinta}}
  {\thesection}{0.5em}{}
  [\vspace{-0.4em}{\color{verdeTerm}\rule{2.5cm}{2pt}}]

\titleformat{\subsection}
  {\large\bfseries\color{verdeOsc}}
  {\thesubsection}{0.5em}{}

\titleformat{\subsubsection}
  {\normalsize\bfseries\color{tinta}}
  {\thesubsubsection}{0.5em}{}

% ---- Encabezado y pie ----
\pagestyle{fancy}
\fancyhf{}
\fancyhead[L]{\small\ttfamily\color{verdeTerm}\textbf{>}\color{tinta2}\ data\textcolor{verdeTerm}{.}lung}
\fancyhead[R]{\small\color{tinta2}Gu\'ia de estudio $\cdot$ defensa TFM}
\fancyfoot[C]{\small\color{mute}\thepage}
\renewcommand{\headrulewidth}{0.4pt}
\renewcommand{\footrulewidth}{0pt}
\setlength{\headheight}{15pt}

% ---- Cajas ----
\newtcolorbox{guioncaja}[1]{
  enhanced, breakable,
  colback=plano, colframe=verdeTerm,
  boxrule=1pt, arc=3pt,
  left=8pt, right=8pt, top=6pt, bottom=6pt,
  title={\bfseries\color{white}#1},
  colbacktitle=verdeTerm,
  attach boxed title to top left={xshift=6pt, yshift=-6pt},
  boxed title style={colback=verdeTerm, sharp corners, boxrule=0pt}
}

\newtcolorbox{tipbox}{
  enhanced, breakable,
  colback=planoOsc, colframe=ambar,
  boxrule=0.8pt, arc=2pt,
  left=8pt, right=8pt, top=5pt, bottom=5pt
}

\newtcolorbox{prcaja}{
  enhanced, breakable,
  colback=white, colframe=linea,
  boxrule=0.6pt, arc=2pt,
  left=8pt, right=8pt, top=4pt, bottom=4pt
}

\newcommand{\gene}[1]{\textit{#1}}
\newcommand{\kicker}[1]{{\ttfamily\color{verdeTerm}\textbf{>} #1}}
\newcommand{\slidebadge}[1]{{\colorbox{verdeTerm}{\color{white}\bfseries\;#1\;}}}

\begin{document}

% =================== PORTADA =========================
\thispagestyle{empty}
\vspace*{2cm}
\begin{center}
  {\ttfamily\fontsize{48pt}{48pt}\selectfont\bfseries
   data\textcolor{verdeTerm}{.}lung}\\[0.2cm]
  {\itshape\large\color{tinta2}Del dataset a una firma g\'enica replicable}\\[0.8cm]

  \textcolor{verdeTerm}{\rule{3cm}{2pt}}\\[0.6cm]

  {\LARGE\bfseries Gu\'ia de estudio para la defensa}\\[0.4cm]
  {\Large\color{tinta2}Framework transcript\'omico para biomarcadores\\en c\'ancer de pulm\'on}\\[1.5cm]

  \begin{tabular}{rl}
    \color{mute}\textbf{Autor}      & Daniel Tapia D\'iez \\
    \color{mute}\textbf{Tutor}      & Leonardo Dulcetti \\
    \color{mute}\textbf{Universidad}& UAX $\cdot$ M\'aster en Bioinform\'atica \\
    \color{mute}\textbf{Defensa}    & 14 de agosto de 2026 \\
  \end{tabular}\\[2cm]

  \begin{tcolorbox}[colback=plano, colframe=verdeTerm, width=0.85\textwidth,
                    arc=3pt, boxrule=1pt]
    \small\itshape
    Este documento tiene dos partes: (A) conceptos clave explicados en lenguaje llano
    para llegar a la defensa con seguridad; (B) gui\'on diapositiva por diapositiva
    con timing y transiciones. A\~n\'adense preguntas probables del tribunal y un
    checklist de \'ultima hora.
  \end{tcolorbox}

  \vfill
  {\ttfamily\small\color{mute}L\'eelo dos veces enteras antes del d\'ia D. La v\'ispera, s\'olo la Parte B.}
\end{center}
\newpage

\tableofcontents
\newpage

% =================== PARTE A =========================
\section*{Parte A $\cdot$ Conceptos clave}
\addcontentsline{toc}{section}{Parte A $\cdot$ Conceptos clave}
{\color{verdeTerm}\rule{2.5cm}{2pt}}

\subsection{Transcript\'omica y expresi\'on g\'enica}

\begin{itemize}[leftmargin=1.2em, itemsep=2pt]
  \item \textbf{ADN}: es el libro de instrucciones de la c\'elula. Todas las c\'elulas de tu cuerpo tienen el mismo ADN.
  \item \textbf{Genes}: son "frases" concretas dentro del libro. Un gen codifica una prote\'ina o una funci\'on reguladora.
  \item \textbf{Expresi\'on g\'enica}: no todos los genes se "leen" al mismo tiempo. Una c\'elula del pulm\'on lee genes distintos que una del h\'igado. La expresi\'on es \emph{cu\'anto} se est\'a leyendo cada gen.
  \item \textbf{Transcriptoma}: el conjunto completo de genes activos en una muestra en un momento dado. "Medir la expresi\'on" es medir el transcriptoma.
  \item \textbf{Por qu\'e importa en c\'ancer}: un tumor tiene un patr\'on distinto del tejido sano, y dentro de los tumores adenocarcinoma y escamoso tienen patrones distintos. Esa \textbf{firma de expresi\'on} sirve para diagn\'ostico y pron\'ostico.
\end{itemize}

\subsection{C\'omo se mide la expresi\'on: microarrays y RNA-Seq}

\begin{itemize}[leftmargin=1.2em, itemsep=2pt]
  \item \textbf{Microarray}: chip de vidrio con miles de sondas de ADN pegadas. Vuelcas ARN de tu muestra, se une a las sondas complementarias y se mide la fluorescencia. Tecnolog\'ia cl\'asica (2000-2015), barata, dominante en estudios GEO antiguos.
  \item \textbf{RNA-Seq}: secuencia directamente todo el ARN. Tecnolog\'ia moderna (2015 en adelante), m\'as precisa, est\'andar en TCGA.
  \item \textbf{Por qu\'e importa aqu\'i}: el framework se entren\'o en \textbf{microarrays} (11 cohortes GEO) y se valid\'o en \textbf{RNA-Seq} (TCGA). Que un panel entrenado en microarray funcione en RNA-Seq demuestra que la se\~nal es real, no un artefacto tecnol\'ogico.
\end{itemize}

\subsection{GEO (Gene Expression Omnibus) y cohortes}

\begin{itemize}[leftmargin=1.2em, itemsep=2pt]
  \item \textbf{GEO}: repositorio p\'ublico del NCBI (EE. UU.) con m\'as de 200.000 estudios de expresi\'on g\'enica.
  \item \textbf{Cohorte / estudio / dataset}: sin\'onimos. Cada uno con identificador tipo \texttt{GSE10245}, \texttt{GSE31210}.
  \item \textbf{Metadatos}: descripciones de cada muestra (edad, tumor/sano, subtipo, estadio\ldots) en \textbf{texto libre}, sin estructura. Por eso el paso 3 del pipeline necesita un LLM: convertir texto no estructurado en etiquetas can\'onicas.
\end{itemize}

\subsection{Adenocarcinoma vs Carcinoma escamoso}

\begin{itemize}[leftmargin=1.2em, itemsep=2pt]
  \item Son los dos subtipos m\'as frecuentes del c\'ancer de pulm\'on \emph{no microc\'itico} (a su vez, el 85\,\% de todos los c\'anceres de pulm\'on).
  \item \textbf{ADC (adenocarcinoma)}: origen glandular, periferia pulmonar. Marcadores cl\'asicos: TTF-1, napsina A.
  \item \textbf{SQC (escamoso)}: origen en el epitelio bronquial, m\'as central. Marcadores cl\'asicos: p40, CK5/6, DSG3.
  \item \textbf{Por qu\'e es cr\'itico distinguirlos}: el tratamiento cambia. \textbf{Pemetrexed} (quimio) y \textbf{bevacizumab} (antiangiog\'enico) est\'an \textbf{contraindicados en escamoso} por riesgo de hemorragia mortal. Un diagn\'ostico err\'oneo puede matar al paciente.
\end{itemize}

\subsection{Inmunohistoqu\'imica (IHC) y panel OMS}

\begin{itemize}[leftmargin=1.2em, itemsep=2pt]
  \item \textbf{IHC}: t\'ecnica del laboratorio de anatom\'ia patol\'ogica. Anticuerpos marcados que se unen a prote\'inas espec\'ificas. Es lo que hace el pat\'ologo al te\~nir y mirar al microscopio.
  \item \textbf{Panel OMS 2015} (Travis et al.): conjunto peque\~no de marcadores recomendado por la OMS para diagn\'ostico rutinario. En el TFM se toma como \textbf{verdad cl\'inica}.
  \begin{itemize}
    \item \textbf{Escamoso}: KRT5, KRT6A, KRT6B, KRT13, KRT14, KRT15, KRT16, KRT17, CDH3, DSG3, PKP1, TP63.
    \item \textbf{Adenocarcinoma}: NKX2-1 (TTF-1), NAPSA (napsina A), SFTPB, SFTPC, SFTPA1, SFTPA2, MUC1, CEACAM5, CEACAM6, KRT7.
  \end{itemize}
  \item \textbf{Validaci\'on del TFM}: la firma de 20 genes seleccionada por LASSO \textbf{recupera 18 de los 20 marcadores del panel OMS} sin que se le dijera nada del panel. Es un aval biol\'ogico fort\'isimo.
\end{itemize}

\subsection{Modelos de lenguaje (LLM), Llama y Groq}

\begin{itemize}[leftmargin=1.2em, itemsep=2pt]
  \item \textbf{LLM}: modelo estad\'istico entrenado con enormes cantidades de texto que aprende a generar y clasificar lenguaje natural.
  \item \textbf{Llama 3.3-70b}: LLM de c\'odigo abierto publicado por Meta. 70.000 millones de par\'ametros. Es el modelo usado para curar los metadatos de GEO.
  \item \textbf{Groq}: infraestructura de inferencia ultrarr\'apida (hardware LPU en vez de GPU). Sirve modelos abiertos como Llama con latencias de decenas de milisegundos, lo que hace viable procesar miles de muestras.
  \item \textbf{Por qu\'e un LLM en bioinform\'atica}: los metadatos GEO est\'an en texto libre, con abreviaturas, faltas y sin\'onimos. Codificar reglas regex para todos los casos es imposible. Un LLM entiende contexto: le pasas ``\emph{primary lung tumor, LUSC, stage III}'' y devuelve la etiqueta can\'onica \texttt{LUSC}.
\end{itemize}

\subsection{Batch effect y por qu\'e integrar cohortes es dif\'icil}

\begin{itemize}[leftmargin=1.2em, itemsep=2pt]
  \item \textbf{Efecto lote (batch effect)}: al integrar muestras de laboratorios distintos, gran parte de la varianza es t\'ecnica (chip, reactivo, operador, a\~no), no biol\'ogica.
  \item \textbf{ComBat}: algoritmo cl\'asico (Johnson et al., 2007) que ``corrige'' el efecto lote normalizando cohortes. Riesgo: si el efecto lote est\'a confundido con la variable de inter\'es, ComBat borra tambi\'en se\~nal biol\'ogica real.
  \item \textbf{En este TFM}: no se aplica ComBat. Se usa \textbf{meta-an\'alisis por consenso}.
\end{itemize}

\subsection{Meta-an\'alisis por consenso}

\begin{itemize}[leftmargin=1.2em, itemsep=2pt]
  \item Cada cohorte se analiza por separado.
  \item Para cada gen y cohorte se calcula el \textbf{tama\~no de efecto d de Cohen}:
  \[ d = \frac{\bar{x}_{\text{tumor}} - \bar{x}_{\text{sano}}}{s_{\text{combinado}}} \]
  \item $|d| > 0{,}5$ es efecto moderado; $|d| > 0{,}8$ grande.
  \item Un gen entra en la firma s\'olo si mantiene \textbf{el mismo signo} y $|d| > 0{,}5$ en al menos las cohortes exigidas (aqu\'i: 3 independientes).
  \item Ventaja frente a ComBat: m\'as conservador, no fuerza normalizaciones. S\'olo sobreviven se\~nales robustas.
\end{itemize}

\subsection{LODO $\cdot$ Leave-One-Dataset-Out}

\begin{itemize}[leftmargin=1.2em, itemsep=2pt]
  \item Variante estricta de validaci\'on cruzada.
  \item Con N cohortes: entrenas con N-1 y eval\'uas en la retirada. Repites N veces.
  \item \textbf{Contraste con CV k-fold cl\'asica}: la CV normal mezcla muestras de todas las cohortes. Dos muestras de la misma cohorte comparten sesgos t\'ecnicos; evaluarlas juntas infla el rendimiento. LODO lo evita.
  \item \textbf{En este TFM}: LODO es el criterio principal. Un AUC = 0,925 en LODO es \textbf{m\'as cre\'ible} que un AUC = 0,95 en k-fold.
\end{itemize}

\subsection{Aprendizaje autom\'atico usado}

\subsubsection*{LASSO (regresi\'on log\'istica con regularizaci\'on L1)}

\begin{itemize}[leftmargin=1.2em, itemsep=2pt]
  \item Modelo lineal con penalizaci\'on que fuerza coeficientes a cero. Selecciona s\'olo los genes \'utiles.
  \item Ventajas: (i) panel peque\~no e interpretable (coeficientes por gen), (ii) determinista con \texttt{random\_state} fijo, (iii) un \'unico hiperpar\'ametro (C).
  \item \textbf{Es el m\'etodo principal del TFM}.
\end{itemize}

\subsubsection*{Random Forest (RF)}

\begin{itemize}[leftmargin=1.2em, itemsep=2pt]
  \item ``Bosque'' de \'arboles de decisi\'on. Cada \'arbol se entrena con submuestra aleatoria y selecci\'on aleatoria de variables; se vota entre \'arboles.
  \item No hace selecci\'on tan agresiva como LASSO, pero da un \textbf{ranking de importancia} \'util como complemento.
  \item En el TFM: AUC 0,933 (ligeramente mejor que LASSO), pero se descarta porque no da un panel interpretable.
\end{itemize}

\subsubsection*{SVM (Support Vector Machine)}

\begin{itemize}[leftmargin=1.2em, itemsep=2pt]
  \item Encuentra el hiperplano \'optimo que separa dos clases en un espacio de alta dimensi\'on. Kernels (RBF) para separaciones no lineales.
  \item En el TFM: AUC 0,957 (el mejor absoluto), pero pesos dif\'iciles de interpretar como ``importancia de gen''.
\end{itemize}

\begin{tipbox}
\textbf{Por qu\'e LASSO si SVM da m\'as AUC?} Interpretabilidad cl\'inica. Diferencia de 3 puntos de AUC no compensa perder un panel que un cl\'inico puede leer.
\end{tipbox}

\subsection{M\'etricas de rendimiento}

\begin{itemize}[leftmargin=1.2em, itemsep=2pt]
  \item \textbf{AUC} (Area Under ROC Curve): probabilidad de que el modelo ordene una muestra positiva por encima de una negativa. Rango 0,5 (azar) a 1,0 (perfecto). Independiente del umbral.
  \item \textbf{Sensibilidad} (recall / TPR): de todos los positivos reales, cu\'antos detecto. Con 100 tumores y 95 detectados, sensibilidad = 0,95.
  \item \textbf{Especificidad}: de todos los negativos reales, cu\'antos descarto correctamente.
  \item \textbf{Balanced Accuracy}: media de sensibilidad y especificidad. \'Util con clases desbalanceadas.
  \item \textbf{PPV / VPP}: valor predictivo positivo. De los que predije positivos, cu\'antos lo son de verdad. Depende de la prevalencia.
\end{itemize}

\begin{tipbox}
\textbf{Regla mnemot\'ecnica.}\\
Sensibilidad: ``no se me escape ning\'un enfermo''.\\
Especificidad: ``no marque como enfermo a alguien sano''.
\end{tipbox}

\subsection{Calibraci\'on (isot\'onica y Platt)}

\begin{itemize}[leftmargin=1.2em, itemsep=2pt]
  \item Un modelo puede tener buen AUC pero puntuaciones \textbf{descalibradas} entre cohortes: predice 0,8 y la probabilidad real es 0,5.
  \item \textbf{Platt scaling}: regresi\'on log\'istica sobre las puntuaciones. \textbf{Isot\'onica}: funci\'on mon\'otona no param\'etrica, m\'as flexible.
  \item \textbf{En el TFM}: la calibraci\'on isot\'onica sube BalAcc de 0,772 $\to$ 0,810 y AUC de 0,925 $\to$ 0,953. Especialmente importante en tumor-vs-sano porque el umbral \'optimo cambia entre cohortes.
\end{itemize}

\subsection{Bootstrap e intervalos de confianza}

\begin{itemize}[leftmargin=1.2em, itemsep=2pt]
  \item \textbf{Bootstrap}: t\'ecnica de remuestreo. Del conjunto de muestras tomas N con reemplazo, recalculas la m\'etrica, repites (p. ej. 1000 veces). La distribuci\'on de esas 1000 m\'etricas te da la incertidumbre.
  \item \textbf{IC 95\,\% bootstrap}: percentiles 2,5 y 97,5 de esa distribuci\'on.
  \item \textbf{TFM}: AUC TCGA = 0,857, IC 95\,\% = [0,812; 0,899]. Si repiti\'eramos el experimento, en el 95\,\% de las repeticiones el AUC caer\'ia entre esos dos valores.
\end{itemize}

\subsection{TCGA y cBioPortal}

\begin{itemize}[leftmargin=1.2em, itemsep=2pt]
  \item \textbf{TCGA} (The Cancer Genome Atlas): consorcio del NIH que caracteriz\'o miles de tumores humanos. Pulm\'on: \textbf{TCGA-LUAD} (adenocarcinoma) y \textbf{TCGA-LUSC} (escamoso).
  \item \textbf{cBioPortal}: portal y API p\'ublica que permite descargar los datos de TCGA limpios. En el TFM se us\'o su REST API.
  \item \textbf{Papel en el TFM}: validaci\'on externa RNA-Seq. Panel LASSO entrenado en microarray aplicado a TCGA (n=408) $\to$ AUC = 0,857. Prueba de que la firma es transferible entre plataformas.
\end{itemize}

\subsection{Correlaci\'on de Spearman ($\rho$)}

\begin{itemize}[leftmargin=1.2em, itemsep=2pt]
  \item Medida de asociaci\'on entre dos variables ordinales.
  \item $\rho = 1$: mon\'otona positiva perfecta; $\rho = -1$: mon\'otona negativa perfecta; $\rho \approx 0$: sin asociaci\'on.
  \item \textbf{En el TFM}: la firma correlaciona negativamente con la composici\'on alv\'eolo-capilar ($\rho = -0{,}626$ media, hasta $\rho = -0{,}858$ en GSE31210). Interpretaci\'on biol\'ogica: cuanto m\'as tumor, menos alv\'eolos.
\end{itemize}

\subsection{Streamlit y la interfaz data.lung}

\begin{itemize}[leftmargin=1.2em, itemsep=2pt]
  \item \textbf{Streamlit}: framework Python para apps web interactivas sin escribir HTML/JS.
  \item \textbf{data.lung}: app multip\'agina del TFM. Portada, framework, resultados (con buscador de gen), cohortes individuales, asistente conversacional Llama acotado.
  \item \textbf{Reproducibilidad}: todas las cifras se leen de los CSV generados por el pipeline. Ning\'un n\'umero est\'a hardcodeado.
\end{itemize}

\subsection{Ejes biol\'ogicos reales (interpretaci\'on)}

La firma no es una caja negra. Al analizarla emergen dos ejes claros:

\begin{itemize}[leftmargin=1.2em, itemsep=2pt]
  \item \textbf{Eje 1 $\cdot$ P\'erdida alv\'eolo-capilar}: genes como SFTPC, SFTPB, AGER, CLDN18 (t\'ipicos de neumocitos y capilares). Cuanto m\'as tumor, menos estos genes.
  \item \textbf{Eje 2 $\cdot$ Actividad proliferativa}: MKI67 (marcador cl\'asico Ki-67), TOP2A, CCNB1, BIRC5, AURKA. Cuanto m\'as agresivo, m\'as proliferaci\'on.
\end{itemize}

Ambos se recuperan \textbf{sin haberlos declarado}: el pipeline los descubre por estad\'istica.

\subsection{NanoString y RT-qPCR (transferibilidad cl\'inica)}

\begin{itemize}[leftmargin=1.2em, itemsep=2pt]
  \item \textbf{NanoString nCounter}: mide expresi\'on de un panel peque\~no de genes sobre tejido fijado en formalina (FFPE). Perfecta para cl\'inica.
  \item \textbf{RT-qPCR}: PCR cuantitativa. Est\'andar cl\'inico para paneles muy peque\~nos.
  \item \textbf{Relevancia}: los 20 genes del panel son medibles por estas t\'ecnicas. El panel podr\'ia trasladarse a un test cl\'inico. Diferencia clave respecto a firmas que requieren microarray/RNA-Seq.
\end{itemize}

\subsection{Reproducibilidad}

\begin{itemize}[leftmargin=1.2em, itemsep=2pt]
  \item \textbf{Definici\'on fuerte}: cualquiera con acceso al c\'odigo y datos puede regenerar todas las cifras y figuras.
  \item \textbf{En este TFM}: pipeline con \texttt{random\_state=42} en todos los pasos aleatorizados. CSV intermedios almacenados. C\'odigo p\'ublico en GitHub.
\end{itemize}

\subsection{Objetivo del TFM en una frase}

\begin{tcolorbox}[colback=plano, colframe=verdeTerm, arc=3pt, boxrule=1pt]
\itshape\small
Framework transcript\'omico reproducible que integra cohortes GEO de c\'ancer de pulm\'on mediante curaci\'on con LLM, meta-an\'alisis por consenso y clasificaci\'on LASSO validada en LODO, con una firma g\'enica interpretable que recupera el panel diagn\'ostico OMS y se transfiere a RNA-Seq (TCGA).
\end{tcolorbox}

Si te preguntan ``en una frase, qu\'e has hecho'', esa es tu respuesta.

\newpage

% =================== PARTE B =========================
\section*{Parte B $\cdot$ Gui\'on diapositiva por diapositiva}
\addcontentsline{toc}{section}{Parte B $\cdot$ Gui\'on por diapositiva}
{\color{verdeTerm}\rule{2.5cm}{2pt}}

\medskip

\textbf{Formato de cada slide.} Timing objetivo, texto sugerido (aprender la intenci\'on, no memorizar palabra por palabra), puntos que enfatizar y transici\'on.

% ----- SLIDE 1 -----
\subsection*{\slidebadge{Slide 1} Portada $\cdot$ 30 s}
\addcontentsline{toc}{subsection}{Slide 1 $\cdot$ Portada}

\begin{guioncaja}{Gui\'on}
Buenos d\'ias. Muchas gracias por estar aqu\'i. Mi nombre es Daniel Tapia D\'iez y presento hoy mi Trabajo Fin de M\'aster, titulado \emph{Framework transcript\'omico basado en la integraci\'on de modelos de lenguaje y aprendizaje autom\'atico para la identificaci\'on de biomarcadores en c\'ancer de pulm\'on}, tutorizado por Leonardo Dulcetti. El trabajo se materializa en un framework al que hemos llamado \textbf{data.lung}, construido sobre software libre --- Python, scikit-learn, Streamlit, Llama sobre Groq y GitHub --- y datos p\'ublicos de NCBI GEO y del portal cBioPortal para TCGA.
\end{guioncaja}

\textbf{Enfatizar}: tu nombre y el nombre del framework (data.lung). ``Software libre y datos p\'ublicos'' = reproducibilidad.\\
\textbf{Transici\'on}: \emph{``Empecemos por el porqu\'e del trabajo.''}

% ----- SLIDE 2 -----
\subsection*{\slidebadge{Slide 2} Problema y objetivo $\cdot$ 45 s}
\addcontentsline{toc}{subsection}{Slide 2 $\cdot$ Problema y objetivo}

\begin{guioncaja}{Gui\'on}
El c\'ancer de pulm\'on sigue siendo la principal causa de muerte por c\'ancer en el mundo: \textbf{2,48 millones de casos anuales} y supervivencia a 5 a\~nos \textbf{inferior al 20\,\%}. Dentro del c\'ancer de pulm\'on no microc\'itico, distinguir adenocarcinoma de escamoso es cr\'itico, porque tratamientos como pemetrexed o bevacizumab est\'an \textbf{contraindicados en escamoso} por riesgo de hemorragia mortal.

Existen cientos de estudios de expresi\'on g\'enica p\'ublicos en GEO, pero integrarlos es dif\'icil: plataformas heterog\'eneas, metadatos en texto libre y firmas g\'enicas publicadas que muchas veces \textbf{no replican} en cohortes independientes. Nuestro objetivo ha sido construir un framework reproducible que combine tres piezas: curaci\'on de metadatos con LLM, validaci\'on externa Leave-One-Dataset-Out y una firma consenso interpretable.
\end{guioncaja}

\textbf{Enfatizar}: cifra cl\'inica (2,48\,M y $<$20\,\%) para dimensionar; palabra \textbf{contraindicados}; palabra \textbf{no replican} para situar el hueco.\\
\textbf{Transici\'on}: \emph{``Veamos c\'omo se ha construido el pipeline.''}

% ----- SLIDE 3 -----
\subsection*{\slidebadge{Slide 3} Pipeline $\cdot$ 1 min 30 s}
\addcontentsline{toc}{subsection}{Slide 3 $\cdot$ Pipeline}

\begin{guioncaja}{Gui\'on}
El pipeline tiene nueve pasos que van del \textbf{dataset crudo} hasta la \textbf{firma interpretable}. Los pasos 1 y 2 hacen la ingesta y limpieza t\'ecnica desde GEO. El paso 3 es el primero clave: usamos \textbf{Llama 3.3 de 70.000 millones de par\'ametros, servido con Groq}, para curar los metadatos cl\'inicos. GEO tiene los metadatos en texto libre y las expresiones regulares se quedan cortas; el LLM entiende contexto. Alcanzamos una tasa de \'exito del \textbf{85,9\,\%}.

Los pasos 4 a 6 son an\'alisis diferencial cohorte a cohorte, calculando \textbf{d de Cohen}. El paso 7 es la \textbf{validaci\'on externa LODO}: cada cohorte se retira una vez y se usa como test independiente. M\'as conservador que la validaci\'on cruzada est\'andar porque respeta la agrupaci\'on por cohorte.

El paso 8 es el \textbf{consenso multi-cohorte}: un gen entra en la firma s\'olo si mantiene el mismo signo y una magnitud m\'inima en las tres cohortes. El paso 9 entrega la firma final y el panel m\'inimo.

Detalle metodol\'ogico importante: no aplicamos ComBat. LODO no corrige el efecto lote, \textbf{lo mide}.
\end{guioncaja}

\textbf{Enfatizar}: 85,9\,\% LLM. La frase ``LODO no corrige el efecto lote, lo mide'' es muy potente --- memor\'izala.\\
\textbf{Transici\'on}: \emph{``Con este pipeline, estos son los resultados.''}

% ----- SLIDE 4 -----
\subsection*{\slidebadge{Slide 4} Firma y panel m\'inimo $\cdot$ 1 min}
\addcontentsline{toc}{subsection}{Slide 4 $\cdot$ Firma y panel}

\begin{guioncaja}{Gui\'on}
Partimos de \textbf{11 cohortes GEO}. Ocho cumplieron los criterios de tama\~no y balance, sumando \textbf{1.157 muestras} curadas por el LLM. El an\'alisis de consenso deja \textbf{1.174 genes validados} en las tres cohortes independientes.

A partir de ese panel amplio, aplicamos LASSO para obtener un \textbf{panel m\'inimo de 20 genes}. El panel de 20 genes obtiene un AUC de \textbf{0,966}, pr\'acticamente indistinguible del 0,970 de la firma completa. Se pierde muy poca informaci\'on y se gana much\'isima aplicabilidad cl\'inica.

En el top 5 aparecen \textbf{DSG3, KRT5, CALML3, KRT6B y PKP1} --- todos genes de desmosomas y queratinizaci\'on, biolog\'ia del epitelio escamoso.
\end{guioncaja}

\textbf{Enfatizar}: los tres bignums (11, 8/11, 1.157) dichos despacio. ``20 genes $\approx$ firma completa'' es la clave cl\'inica.\\
\textbf{Transici\'on}: \emph{``\'{}Y estos genes son biol\'ogicamente cre\'ibles? Miremos la validaci\'on externa.''}

% ----- SLIDE 5 -----
\subsection*{\slidebadge{Slide 5} IHC y biolog\'ia $\cdot$ 1 min 30 s}
\addcontentsline{toc}{subsection}{Slide 5 $\cdot$ IHC y biolog\'ia}

\begin{guioncaja}{Gui\'on}
La OMS publica desde 2015 un panel de marcadores IHC para el diagn\'ostico rutinario de c\'ancer de pulm\'on: 12 marcadores de escamoso y 8 de adenocarcinoma. La firma obtenida, seleccionada de forma \textbf{puramente estad\'istica sin conocer ese panel}, recupera \textbf{12 de 12 marcadores escamosos y 6 de 8 adenocarcinomatosos: 18 de 20 en total, un 90\,\%}. Coincidencia entre selecci\'on estad\'istica y pr\'actica cl\'inica: validaci\'on externa muy fuerte.

Analizando la estructura interna emergen dos ejes biol\'ogicos claros. El \textbf{primer eje} captura la p\'erdida alv\'eolo-capilar: correlaci\'on negativa con marcadores alveolares con $\rho$ media de $-0{,}626$, llegando a $\rho = -0{,}858$ en GSE31210, con p-valor del orden de $10^{-66}$. El \textbf{segundo eje} captura la actividad proliferativa: MKI67, TOP2A, CCNB1, BIRC5, AURKA --- marcadores cl\'asicos de proliferaci\'on tumoral.

Es decir, la firma no es una caja negra. Captura la transici\'on biol\'ogica del pulm\'on sano al tejido tumoral en dos ejes que un pat\'ologo reconoce.
\end{guioncaja}

\textbf{Enfatizar}: \textbf{18/20 sin haberlo declarado} es tu argumento m\'as fuerte. $\rho = -0{,}858$, $p \approx 10^{-66}$ son cifras impresionantes.\\
\textbf{Transici\'on}: \emph{``En cuanto a rendimiento, hemos ido m\'as all\'a del AUC b\'asico.''}

% ----- SLIDE 6 -----
\subsection*{\slidebadge{Slide 6} Rendimiento, recalibraci\'on y TCGA $\cdot$ 1 min 30 s}
\addcontentsline{toc}{subsection}{Slide 6 $\cdot$ Rendimiento y TCGA}

\begin{guioncaja}{Gui\'on}
En la tarea \textbf{tumor vs sano}, el LODO base da un AUC de 0,925 y una BalAcc de 0,772. Al aplicar \textbf{recalibraci\'on isot\'onica}, el AUC sube a \textbf{0,953} y la BalAcc a \textbf{0,810}. La especificidad mejora especialmente, de 0,561 a 0,632, reduciendo falsos positivos.

Comparamos con Random Forest y SVM: LASSO 0,925, Random Forest 0,933, SVM 0,957. SVM gana ligeramente, pero elegimos LASSO por interpretabilidad --- un panel de 20 genes con coeficientes expl\'icitos es m\'as \'util cl\'inicamente que un hiperplano SVM.

Validaci\'on en \textbf{RNA-Seq}: aplicamos el panel LASSO entrenado en microarray a \textbf{TCGA-LUAD y TCGA-LUSC}, 408 muestras. Ninguna muestra TCGA particip\'o en la selecci\'on ni en la calibraci\'on. AUC = \textbf{0,857, con IC 95\,\% [0,812; 0,899] por bootstrap}. La firma se transfiere a una plataforma tecnol\'ogica distinta, lo que refuerza que la se\~nal es biol\'ogica y no un artefacto de microarray.
\end{guioncaja}

\textbf{Enfatizar}: BalAcc 0,772 $\to$ 0,810 tras recalibraci\'on. AUC 0,857 en TCGA como sello de garant\'ia externa. ``Ninguna muestra TCGA particip\'o'' cierra dudas sobre \emph{data leakage}.\\
\textbf{Transici\'on}: \emph{``Y ahora, la parte m\'as divertida. Voy a ense\~naros la herramienta.''}

% ----- SLIDE 7 -----
\subsection*{\slidebadge{Slide 7} data.lung $\cdot$ DEMO EN VIVO $\cdot$ 2 min}
\addcontentsline{toc}{subsection}{Slide 7 $\cdot$ Demo}

\begin{guioncaja}{Gui\'on ANTES de cambiar al navegador}
Todo lo que os he contado se recoge en una herramienta interactiva multip\'agina que hemos llamado \textbf{data.lung}, hecha en Streamlit. Un buscador de gen en tiempo real, vistas por cohorte, y un asistente conversacional basado en Llama acotado a los resultados del trabajo --- no responde nada que no est\'e en los datos.

Voy a ense\~n\'aroslo un momento.
\end{guioncaja}

\textbf{Cambia a la app} (Cmd+Tab). \textbf{Secuencia recomendada}:

\begin{enumerate}[leftmargin=1.4em, itemsep=3pt]
  \item \textbf{Portada} (10 s): ``Esta es la portada.''
  \item \textbf{Framework} (15 s): ``Aqu\'i las cifras del framework. Todo se lee de CSV, ninguna cifra est\'a en el c\'odigo.''
  \item \textbf{Resultados $\to$ Buscador de gen} (45 s):
    \begin{itemize}
      \item Escribe \texttt{DSG3}: ``DSG3 aparece: es uno de los cinco genes principales del panel escamoso.''
      \item Escribe \texttt{SFTPC}: ``SFTPC no aparece en el panel m\'inimo, aunque s\'i en la firma completa. Es un ejemplo de c\'omo LASSO reduce redundancia.''
    \end{itemize}
  \item \textbf{Asistente} (30 s):
    \begin{itemize}
      \item Pregunta: \emph{``qu\'e biomarcadores identifica el framework?''}
      \item Coment\'ario: ``Fijaos que responde con los datos concretos del trabajo, y si le pregunto algo que no est\'a en los resultados, me lo dice expl\'icitamente.''
    \end{itemize}
\end{enumerate}

\textbf{Cambia de vuelta al PDF} (Cmd+Tab): ``Volvemos al PDF.''

\begin{tipbox}
\textbf{Antes de empezar la defensa}, tener corriendo:\\[3pt]
\texttt{streamlit run agentes/paso12\_web\_chatbot.py}\\[3pt]
Y la app abierta en el navegador, minimizada. Cmd+Tab probado.
\end{tipbox}

\textbf{Transici\'on}: \emph{``Recapitulemos las conclusiones.''}

% ----- SLIDE 8 -----
\subsection*{\slidebadge{Slide 8} Conclusiones $\cdot$ 45 s}
\addcontentsline{toc}{subsection}{Slide 8 $\cdot$ Conclusiones}

\begin{guioncaja}{Gui\'on}
En resumen: un LLM cura los metadatos GEO con \textbf{85,9\,\% de \'exito}, escalando la integraci\'on multi-cohorte sin intervenci\'on manual. El consenso multi-cohorte m\'as LODO produce una firma robusta de \textbf{1.174 genes con AUC 0,968} en subtipo. El panel m\'inimo de \textbf{20 genes es pr\'acticamente equivalente} a la firma completa. La firma \textbf{recupera 18 de 20 marcadores IHC cl\'inicos sin declararlos}. La recalibraci\'on isot\'onica mejora significativamente el rendimiento tumor-vs-sano. Y la firma \textbf{se transfiere a RNA-Seq} con AUC 0,857 en TCGA. Todo el framework es reproducible extremo a extremo, con interfaz web p\'ublica y asistente conversacional acotado.
\end{guioncaja}

\textbf{Enfatizar}: es la \'ultima oportunidad de meter las cifras clave. Dilas.\\
\textbf{Transici\'on}: \emph{``Y de aqu\'i, \'{}por d\'onde continuar\'ia el trabajo?''}

% ----- SLIDE 9 -----
\subsection*{\slidebadge{Slide 9} Trabajo futuro y stack $\cdot$ 45 s}
\addcontentsline{toc}{subsection}{Slide 9 $\cdot$ Trabajo futuro}

\begin{guioncaja}{Gui\'on}
El siguiente paso natural es una \textbf{validaci\'on cl\'inica prospectiva} sobre biopsias frescas con NanoString o RT-qPCR --- el panel de 20 genes est\'a pensado para ser medible por esas t\'ecnicas. Otras l\'ineas: ampliar a subtipos neuroendocrinos, sustituir el LLM externo por un modelo local por privacidad y coste, extender la metodolog\'ia a otros tumores como mama, colon o hepatocarcinoma, y explorar firmas pron\'osticas con regresi\'on de Cox sobre TCGA.

Todo el trabajo est\'a construido sobre software libre y datos p\'ublicos, y disponible en GitHub.
\end{guioncaja}

\textbf{Enfatizar}: \textbf{validaci\'on prospectiva} = camino hacia cl\'inica. Software libre + datos p\'ublicos + GitHub = \'etica de investigaci\'on abierta.\\
\textbf{Transici\'on}: \emph{``Y con esto termino.''}

% ----- SLIDE 10 -----
\subsection*{\slidebadge{Slide 10} Gracias $\cdot$ 15 s}
\addcontentsline{toc}{subsection}{Slide 10 $\cdot$ Cierre}

\begin{guioncaja}{Gui\'on}
Muchas gracias por vuestra atenci\'on. Estoy a vuestra disposici\'on para las preguntas que quer\'ais plantearme.
\end{guioncaja}

Sonr\'ie, respira, silencio corto. Espera el turno de preguntas.

\newpage

% =================== PARTE C =========================
\section*{Parte C $\cdot$ Preguntas probables del tribunal}
\addcontentsline{toc}{section}{Parte C $\cdot$ Preguntas del tribunal}
{\color{verdeTerm}\rule{2.5cm}{2pt}}

\medskip

\begin{prcaja}
\textbf{Q1. \'{}Por qu\'e no ha usado ComBat para corregir efecto lote?}\\[4pt]
\textbf{R.} Porque en varias cohortes GEO el efecto lote est\'a confundido con la variable de inter\'es --- hay estudios que s\'olo contienen tumores, o s\'olo un subtipo. En ese escenario ComBat podr\'ia borrar se\~nal biol\'ogica genuina. Prefiero un enfoque m\'as conservador: normalizar cada cohorte por separado y exigir consenso en varias cohortes independientes. LODO complementa esto midiendo el rendimiento en la peor situaci\'on posible: una cohorte totalmente no vista.
\end{prcaja}

\begin{prcaja}
\textbf{Q2. \'{}C\'omo garantiza que no hay data leakage entre el descubrimiento de la firma y la validaci\'on TCGA?}\\[4pt]
\textbf{R.} Toda la selecci\'on de genes y la calibraci\'on se hizo exclusivamente sobre cohortes GEO en microarray. Las 408 muestras de TCGA se usaron s\'olo en la fase final, aplicando el modelo LASSO ya entrenado. Ninguna muestra TCGA particip\'o en el ajuste ni en la selecci\'on de hiperpar\'ametros.
\end{prcaja}

\begin{prcaja}
\textbf{Q3. El LLM tiene sesgos. \'{}No los introduce en los datos?}\\[4pt]
\textbf{R.} S\'i, es una preocupaci\'on real. Por eso el LLM se restringe a una tarea muy delimitada: leer una descripci\'on de muestra y devolver una etiqueta can\'onica de grupo experimental. No participa en el an\'alisis estad\'istico ni en la selecci\'on de genes. Adem\'as, la tasa de \'exito del 85,9\,\% se audit\'o manualmente sobre un subconjunto para descartar sesgos sistem\'aticos.
\end{prcaja}

\begin{prcaja}
\textbf{Q4. \'{}Por qu\'e LASSO si SVM da un AUC mejor?}\\[4pt]
\textbf{R.} Porque un panel de 20 genes con coeficientes expl\'icitos por gen es mucho m\'as \'util cl\'inicamente que un hiperplano SVM. La diferencia de 3 puntos de AUC no compensa perder la interpretabilidad, sobre todo cuando el objetivo es un panel medible por NanoString o RT-qPCR.
\end{prcaja}

\begin{prcaja}
\textbf{Q5. \'{}La firma es progn\'ostica o s\'olo diagn\'ostica?}\\[4pt]
\textbf{R.} En este trabajo, \'unicamente diagn\'ostica: distinguir ADC de SQC y tumor de sano. La firma pron\'ostica requerir\'ia integrar tiempos de supervivencia y regresi\'on de Cox, y est\'a se\~nalada como l\'inea de trabajo futuro sobre TCGA.
\end{prcaja}

\begin{prcaja}
\textbf{Q6. \'{}Qu\'e pasa si un gen del panel no funciona por RT-qPCR en un laboratorio concreto?}\\[4pt]
\textbf{R.} El panel es redundante: dentro de los 20 genes hay grupos de marcadores del mismo linaje (varias queratinas y desmoglein\'as para escamoso; varios surfactantes para adenocarcinoma). Si uno falla t\'ecnicamente, otros del mismo grupo mantienen la se\~nal.
\end{prcaja}

\begin{prcaja}
\textbf{Q7. \'{}Por qu\'e usar cohortes tan antiguas de microarray si RNA-Seq es el est\'andar hoy?}\\[4pt]
\textbf{R.} Porque hay m\'as cohortes de microarray con metadatos ricos para adenocarcinoma vs escamoso, y eso permite el consenso multi-cohorte. RNA-Seq entra como validaci\'on externa para confirmar transferibilidad, no como base del descubrimiento.
\end{prcaja}

\begin{prcaja}
\textbf{Q8. \'{}No hay riesgo de que el LLM ``invente'' etiquetas?}\\[4pt]
\textbf{R.} El prompt est\'a muy acotado: se le dan las descripciones de las muestras y una lista cerrada de etiquetas posibles, con instrucci\'on expl\'icita de devolver \texttt{unknown} si no encaja. La tasa de \'exito del 85,9\,\% se calcula sobre auditor\'ia manual, y el 14,1\,\% restante se descarta, no se fuerza.
\end{prcaja}

\begin{prcaja}
\textbf{Q9. \'{}Qu\'e le distingue de firmas comerciales como Pervenio Lung RS u Oncotype DX Lung?}\\[4pt]
\textbf{R.} Tres cosas: (i) validaci\'on LODO estricta en 3 cohortes totalmente independientes, no partici\'on aleatoria; (ii) verificaci\'on cruzada contra el panel diagn\'ostico OMS, que Pervenio no reporta; (iii) reproducibilidad completa con c\'odigo y datos p\'ublicos. Los ensayos comerciales son cajas negras protegidas por patente.
\end{prcaja}

\begin{prcaja}
\textbf{Q10. \'{}C\'omo escalar\'ia data.lung a otros tumores?}\\[4pt]
\textbf{R.} El framework es gen\'erico: cambiar la fuente de datos (otras cohortes GEO y su equivalente TCGA), reentrenar la curaci\'on del LLM con las etiquetas del nuevo tumor y ejecutar el mismo pipeline. La l\'ogica de consenso, LODO y selecci\'on LASSO no depende del tejido.
\end{prcaja}

\newpage

% =================== PARTE D =========================
\section*{Parte D $\cdot$ Checklist v\'ispera de la defensa}
\addcontentsline{toc}{section}{Parte D $\cdot$ Checklist}
{\color{verdeTerm}\rule{2.5cm}{2pt}}

\medskip

\begin{itemize}[leftmargin=1.2em, itemsep=5pt, label={$\square$}]
  \item Port\'atil cargado + cargador.
  \item Adaptador HDMI/USB-C si el aula lo requiere.
  \item PDF de la presentaci\'on en local \textbf{y} en pen drive \textbf{y} en Drive/Dropbox.
  \item Streamlit corriendo antes de entrar al aula:\\ \texttt{streamlit run agentes/paso12\_web\_chatbot.py}
  \item Navegador con data.lung abierto y minimizado (Cmd+Tab probado).
  \item Copia impresa opcional (por si falla todo).
  \item Botella de agua.
  \item Reloj visible (o el m\'ovil en modo avi\'on encima de la mesa).
  \item Respirar antes de empezar. Tres segundos de silencio son mucho menos largos de lo que crees.
\end{itemize}

\vspace{2cm}

\begin{tcolorbox}[colback=plano, colframe=verdeTerm, arc=3pt, boxrule=1pt]
\centering\itshape
Suerte, Daniel. La memoria es rigurosa, la app funciona, los datos hablan por s\'i solos.\\
Cuenta la historia con la tranquilidad de saber que todo est\'a donde tiene que estar.
\end{tcolorbox}

\end{document}
